\chapter{Аудит связывающих ядер композитного
\texorpdfstring{$SU(2)$}{SU(2)}-синглета}
\label{ch:v10-discrete-resolution-su2-singlet-binding-audit}

Предыдущий гейт построил активный синглетный проектор $P_S$, но условно
записал родитель $I-P_S$. Теперь проверяется, следует ли такое разделение
каналов из унаследованной динамики.

Коммутант диагонального действия $SU(2)$ на $\mathbf2\otimes\mathbf2$
двумерен. Поэтому всякое инвариантное эрмитово ядро имеет вид
\begin{equation}
 K_{\rm inv}=aP_S+bP_T,\qquad P_T=I-P_S.
\end{equation}
Симметрия фиксирует два канала, но не выбирает разность энергий
$\Delta_{\rm ST}=b-a$, её знак или нормировку.

Канонический обменный оператор существует:
\begin{equation}
 X=\sum_{r=1}^{3}T_r\otimes T_r
   =\frac14I-P_S,
 \qquad
 \Spec X=\left\{-\frac34^{(1)},\frac14^{(3)}\right\}.
\end{equation}
После ручной нормировки он точно воспроизводит условный родитель,
\begin{equation}
 \frac34I+X=P_T.
\end{equation}
Тем самым форма привлекательного синглетного канала канонична. Однако
множитель перед $X$ не порождён действием, а flavor-слой остаётся
тождественным оператором ранга $256$.

Проверены одиннадцать кандидатов: унаследованное нулевое ядро, скалярный
контакт, оператор $T_1\!\cdot T_2$, антисимметризатор, условный $P_T$,
обмен gauge-бозоном, flavor-четырёхфермионный канал, обмен через ванну,
спектральная самоэнергия, portal splitting и fitted composite pole.
Применены шесть критериев: $SU(2)$-тип, присутствие взаимодействия,
синглетное притяжение, родитель коэффициента, выбор flavor-линии и
некруговой масштаб. Критериальная матрица $11\times6$ имеет ранг $6$,
оценки
\begin{equation}
 (4,3,3,3,2,3,4,2,4,3,3)
\end{equation}
и ни одного полного прохода.

\begin{theorem}[Форма обмена выведена, сила связывания --- нет]
Диагональная $SU(2)$-симметрия оставляет ровно два операторных канала, а
$\sum_rT_r\otimes T_r$ канонически различает синглет и триплет. Но она не
определяет ненулевой коэффициент $\Delta_{\rm ST}$, flavor-проектор или
масштаб композитного полюса.
\end{theorem}

\begin{nogocriterion}[Нормировка проектора не является динамикой]
Равенство $3I/4+X=P_T$ доказывает допустимую архитектуру, но не разрешает
назначить её гамильтонианом с единичным коэффициентом. Требуется
унаследованный vertex или общий родитель, независимо выбирающий знак и
величину обменного коэффициента.
\end{nogocriterion}

\begin{protocol}\begin{enumerate}
\item каноническая операторная архитектура: \textbf{$6/6$};
\item охват кандидатов: \textbf{$11/11$};
\item полные проходы: \textbf{$0/11$};
\item физическое происхождение: \textbf{$0/3$};
\item следующий гейт: \textbf{родитель коэффициента связывания}.
\end{enumerate}\end{protocol}

Машинный сертификат сохранён в
\path{s2t_v10_cell_birth_four_volume_nonequilibrium_bath_discrete_resolution_transmuted_mode_asymptotically_free_su2_gauge_singlet_binding_kernel_candidate_audit_gate_results.json}.